% part: intuitionistic-logic
% chapter: propositions-as-types
% section: reduction

\documentclass[../../../include/open-logic-section]{subfiles}

\begin{document}

\olfileid[id]{int}{pty}{red}

\olsection{Reduksi}

Dalam !!{derivation}s deduksi alami, !!a{introduction} yang diikuti
!!a{elimination} merupakan jalan memutar yang dapat ``direduksi''.
Misalnya, \Intro{\land} yang diikuti \Elim{\land} ``saling meniadakan'',
sebab kesimpulan \Elim{\land} selalu merupakan salah satu konjung dari
konjungsi yang dihasilkan \Intro{\land}, sedangkan kedua premis
\Intro{\land} adalah kedua konjung tersebut. Jadi, jika kita memerlukan
!!a{derivation} bagi salah satu konjung, kita cukup menggunakan
!!{derivation} itu tanpa memperkenalkan konjungsi lalu mengeliminasinya.
Kita mempunyai
\begin{gather*}
  \bottomAlignProof
  \AxiomC{}
  \RightLabel{$\delta_1$}
  \DeduceC{$!A_1$}
  \AxiomC{}
  \RightLabel{$\delta_2$}
  \DeduceC{$!A_2$}
  \RightLabel{$\Intro{\land}$}
  \BinaryInfC{$!A_1 \land !A_2$}
  \RightLabel{$\Elim{\land}_i$}
  \UnaryInfC{$!A_i$}
  \DisplayProof\bottomAlignProof
  \quad
  \redone
  \quad
  \AxiomC{}
  \RightLabel{$\delta_i$}
  \DeduceC{$!A_i$}
  \DisplayProof
\end{gather*}

Kita dapat mempertimbangkan relasi reduksi serupa pada term bukti terkait
yang disarankan oleh penyederhanaan ini. Jika $N_1$ merupakan term bukti
bagi~$\delta_1$ dan $N_2$ term bukti bagi~$\delta_2$, kita dapat mengatakan
bahwa term bukti bagi bukti di sebelah kiri tereduksi (dalam satu langkah)
menjadi term bukti bagi bukti di sebelah kanan:
\[
  \proj{i}{\pair{N_1}{N_2}} \redone N_i
\]
Dalam kalkulus lambda bertipe, ini adalah aturan \emph{reduksi-$\beta$}
bagi tipe hasil kali.

Dalam deduksi alami, sepasang inferensi seperti \Intro{\land} yang diikuti
\Elim{\land}, yakni pasangan yang dapat direduksi, disebut
\emph{potongan}. Dalam kalkulus lambda bertipe, term di sebelah kiri
$\redone$ disebut \emph{redex}, sedangkan term di sebelah kanan disebut
\emph{kontraktum}. Berbeda dengan kalkulus lambda tak bertipe, yang hanya
menganggap $(\lambd[x][N])Q$ sebagai redex, kalkulus lambda bertipe
memperluas sintaksnya dengan term yang melibatkan $\pair{N}{M}$,
$\proj{i}{N}$, $\inj{i}{!A}{N}$,
$\dcase{N}{x_1}{M_1}{x_2}{M_2}$, dan $\abort{!A}{N}$, beserta redex yang
bersesuaian.

Demikian pula, kita mempunyai reduksi bagi disjungsi:
\begin{gather*}
  \bottomAlignProof
  \AxiomC{}
  \RightLabel{$\delta$}
  \DeduceC{$!A_i$}
  \RightLabel{$\Intro{\lor}$}
  \UnaryInfC{$!A_1 \lor !A_2$}
  \AxiomC{$\Discharge{!A_1}{x_1}$}
  \RightLabel{$\delta_1$}
  \DeduceC{$!C$}
  \AxiomC{$\Discharge{!A_2}{x_2}$}
  \RightLabel{$\delta_2$}
  \DeduceC{$!C$}
  \DischargeRule{$\Elim{\lor}$}{x_1,x_2}
  \TrinaryInfC{$!C$}
  \DisplayProof\bottomAlignProof
  \quad
  \redone
  \quad
  \AxiomC{}
  \RightLabel{$\delta$}
  \DeduceC{$!A_i$}
  \RightLabel{$\delta_i$}
  \DeduceC{$!C$}
  \DisplayProof
\end{gather*}
Hal ini bersesuaian dengan reduksi pada term bukti:
\begin{gather*}
  \dcase{\inj[!A_i]{i}{M}}{x_1}{N_1}{x_2}{N_2}
  \redone \Subst{N_i}{M}{x_i}
\end{gather*}
dengan $M$ sebagai term bukti bagi~$\delta$, dan $N_i$ sebagai term bukti
bagi~$\delta_i$. Ini adalah aturan reduksi-$\beta$ bagi tipe jumlah. Di
sini, $\Subst{M}{N}{x}$ berarti mengganti semua kemunculan bebas variabel~$x$
dalam $M$ dengan~$N$.

Reduksi tipe fungsi bersesuaian dengan penghapusan jalan memutar berupa
\Intro\lif yang diikuti \Elim\lif.
\begin{gather*}
  \bottomAlignProof
  \AxiomC{$\Discharge{!A}{x}$}
  \RightLabel{$\delta$}
  \DeduceC{$!B$}
  \DischargeRule{$\Intro{\lif}$}{x}
  \UnaryInfC{$!A \lif !B$}
  \AxiomC{}
  \RightLabel{$\delta'$}
  \DeduceC{$!A$}
  \RightLabel{$\Elim{\lif}$}
  \BinaryInfC{$!B$}
  \DisplayProof\bottomAlignProof
  \quad
  \redone
  \quad
  \AxiomC{}
  \RightLabel{$\delta'$}
  \DeduceC{$!A$}
  \RightLabel{$\delta$}
  \DeduceC{$!B$}
  \DisplayProof
\end{gather*}
Bagi term bukti, hal ini sama dengan reduksi-$\beta$ biasa:
\begin{gather*}
  (\lambd[\typeof{x}{!A}][N]M)
  \redone \Subst{N}{M}{x}
\end{gather*}

Selain konversi reduksi pada !!{proof}s, kita dapat mempertimbangkan
konversi permutasi. Konversi ini berlaku pada (sub-)!!{proof}s yang berakhir
dengan \Elim{\lor} lalu diikuti aturan !!{elimination} lain, misalnya:
\begin{multline*}
  \bottomAlignProof
  \AxiomC{}
  \RightLabel{$\delta$}
  \DeduceC{$!A_i$}
  \RightLabel{$\Intro{\lor}$}
  \UnaryInfC{$!A_1 \lor !A_2$}
  \AxiomC{$\Discharge{!A_1}{x_1}$}
  \RightLabel{$\delta_1$}
  \DeduceC{$!C_1 \land !C_2$}
  \AxiomC{$\Discharge{!A_2}{x_2}$}
  \RightLabel{$\delta_2$}
  \DeduceC{$!C_1 \land !C_2$}
  \DischargeRule{$\Elim{\lor}$}{x_1,x_2}
  \TrinaryInfC{$!C_1 \land !C_2$}
  \RightLabel{\Elim{\land}[i]}
  \UnaryInfC{$!C_i$}
  \DisplayProof\bottomAlignProof
  \quad
  \redone
  \\
  \AxiomC{}
  \RightLabel{$\delta$}
  \DeduceC{$!A_i$}
  \RightLabel{$\Intro{\lor}$}
  \UnaryInfC{$!A_1 \lor !A_2$}
  \AxiomC{$\Discharge{!A_1}{x_1}$}
  \RightLabel{$\delta_1$}
  \DeduceC{$!C_1 \land !C_2$}
  \RightLabel{\Elim{\land}[i]}
  \UnaryInfC{$!C_i$}
  \AxiomC{$\Discharge{!A_2}{x_2}$}
  \RightLabel{$\delta_2$}
  \DeduceC{$!C_1 \land !C_2$}
  \RightLabel{\Elim{\land}[i]}
  \UnaryInfC{$!C_i$}
  \DischargeRule{$\Elim{\lor}$}{x_1,x_2}
  \TrinaryInfC{$!C_i$}
  \DisplayProof
\end{multline*}
Jika term bukti bagi $\delta$, $\delta_1$, dan $\delta_2$ masing-masing
adalah $M$, $N_1$, dan $N_2$, aturan reduksi yang bersesuaian bagi term
bukti, yakni term-$\lambd$, adalah:
\[
  \proj{i}{\dcase{M}{x_1}{N_1}{x_2}{N_2}} \redone
  \dcase{M}{x_1}{\proj{i}{N_1}}{x_2}{\proj{i}{N_2}}
\]

Tentu saja, kita dapat menghapus jalan memutar bukan hanya pada akhir
!!{proof}s, melainkan juga di dalam !!{proof}s, dengan menerapkan konversi reduksi
pada sub-!!{proof} yang berakhir dengan jalan memutar. Demikian pula,
reduksi-$\beta$ berlaku pada subterm dari term-$\lambd$. Jika $N$ merupakan
subterm dari $M$ dan $N \redone N'$, kita dapat mengganti subterm $N$
dalam~$M$ dengan $N'$ untuk memperoleh $M'$. Dalam hal ini kita juga menulis
$M \redone M'$. Jika terdapat barisan
$M=M_1 \redone M_2 \redone M_3 \redone \dots \redone M_k=M'$ (termasuk
kasus $k=1$, yakni $M=M'$), kita menulis $M \red M'$.

!!^a{proof} yang tidak memungkinkan penerapan konversi pada sub-!!{proof}
mana pun disebut \emph{normal}. Demikian pula, term-$\lambd$ yang tidak
memungkinkan penerapan konversi (pada subterm mana pun) disebut
\emph{normal}. Term-$\lambd$ normal juga disebut \emph{bentuk normal}.

\end{document}
